% Book pp. 116-122 (PDF pp. 117-123)
\section{Applying the Fundamental Theorem of Algebra}

The \textbf{Fundamental Theorem of Algebra} tells us something very important about
polynomial functions. It says: `Every polynomial function of degree $n$ (where $n >
0$) has at least one complex zero.'

This might sound complicated at first, but here is what it means:
\begin{enumerate}
  \item The degree of a polynomial is the highest power of $x$ in the expression.
  \item For example, $x^3 - 2x^2 + 4x - 8$ has a degree of 3.
  \item Complex numbers include real numbers and imaginary numbers (numbers
  with $i$, where $i = \sqrt{-1}$)
  \item In factorising the polynomial, each root c$_i$ corresponds to a linear factor of
  the form $(x - c_i)$. This means we can write the polynomial as f(x) $= a\,(x -
  c_1)\,(x - c_2) \cdots (x - c_n)$.

  The, $a$ is a non-zero number that affects how the polynomial behaves, and
  $c_1, c_2, \ldots, c_n$, are the roots, or zeros.
\end{enumerate}
Let us go through these examples to show how it works.

\begin{example}{Example 3.11}
Factorise the polynomial f$(x) = 3x^5 - 48x$ completely and find all its zeros.

State the multiplicity of each zero.

\solution
f$(x) = 3x^5 - 48x$

Factorising:

$3x(x^4 - 16)$

$3x[(x^2)^2 - 4^2]$

$3x[(x^2 + 4)(x^2 - 4)]$

$3x[x^2 - (-4)(x^2 - 4)]$

$3x[(x + 2i)(x - 2i)(x + 2)(x - 2)]$

To find the zeros, equate the factors to zero:

$3x[(x + 2i)(x - 2i)(x + 2)(x - 2)] = 0$

$3x = 0$, $x = 0$

$x + 2i = 0$, $x = -2i$

$x - 2i = 0$, $x = 2i$

$x + 2 = 0$, $x = -2$

$x - 2 = 0$, $x = 2$

Therefore, the zeros of f($x$) are 0, 2, -- 2, 2$i$ and -- 2$i$. Since each factor occurs
only once, all the zeros are of multiplicity 1 and the total number of zeros is five.
\end{example}

\begin{example}{Example 3.12}
Factorise the polynomial f$(x) = x^4 - 625$ completely and find all its zeros.

State the multiplicity of each zero.

\solution
f$(x) = x^4 - 625$

Factorising:

$(x^4 - 625)$

$[(x^2)^2 - 25^2]$

$[(x^2 + 25)(x^2 - 25)]$

$3x[x^2 - (-25)(x^2 - 25)]$
\booknote[S3-13]{the ``$3x$'' is carried over from Example 3.11; there is no such
factor here.}

$[(x + 5i)(x - 5i)(x + 5)(x - 5)]$

To find the zeros, equate the factors to zero:

$[(x + 5i)(x - 5i)(x + 5)(x - 5)] = 0$

$x + 5i = 0$, $x = -5i$

$x - 5i = 0$, $x = 5i$

$x + 5 = 0$, $x = -5$

$x - 5 = 0$, $x = 5$

Therefore, the zeros of f($x$) are 5, -- 5, 5$i$ and -- 5$i$. Since each factor occurs
only once, all the zeros are of multiplicity 1 and the total number of zeros is four.
\end{example}

\begin{example}{Example 3.13}
Factorise the polynomial f$(x) = 15x^4 + 8x^2 + 1$ completely and find all its zeros.

State the multiplicity of each zero.

\solution
f$(x) = 15x^4 + 8x^2 + 1$

in this scenario, we can represent $x^2$ by any other letter, for example, u

$f(x) = 15(x^2)^2 + 8x^2 + 1$

$f(x) = 15(u)^2 + 8u + 1$

$15u^2 + 8u + 1$

Factorising:

$15u^2 + 5u + 3u + 1$

$5u(3u + 1) + (3u + 1)$

$(5u + 1)(3u + 1)$

Substituting $x^2$ back,

$(5x^2 + 1)(3x^2 + 1)$

To find the zeros, equate the factors to zero:

$(5x^2 + 1)(3x^2 + 1) = 0$

$5x^2 + 1 = 0$ or $3x^2 + 1 = 0$

For $5x^2 + 1 = 0$

$5x^2 = -1$

$x^2 = -\frac{1}{5}$

$x = \pm\sqrt{-\frac{1}{5}}$

$x = -\sqrt{-\frac{1}{5}}$ or $x = \sqrt{-\frac{1}{5}}$

For $x = -\sqrt{-\frac{1}{5}}$

$x = -\sqrt{\frac{1}{5}} \times \sqrt{-1}$

$x = -i\sqrt{\frac{1}{5}}$

$x = -\frac{\sqrt{5}}{5}\,i$

For $x = \sqrt{-\frac{1}{5}}$

$x = \sqrt{\frac{1}{5}} \times \sqrt{-1}$

$x = i\sqrt{\frac{1}{5}}$

$x = \frac{\sqrt{5}}{5}\,i$

For $3x^2 + 1 = 0$

$3x^2 = -1$

$x^2 = -\frac{1}{3}$

$x = \pm\sqrt{-\frac{1}{3}}$

$x = -\sqrt{-\frac{1}{3}}$ or $x = \sqrt{-\frac{1}{3}}$

For $x = -\sqrt{-\frac{1}{3}}$

$x = -\sqrt{\frac{1}{3}} \times \sqrt{-1}$

$x = -i\sqrt{\frac{1}{3}}$

$x = -\frac{\sqrt{3}}{3}\,i$

For $x = \sqrt{-\frac{1}{3}}$

$x = \sqrt{\frac{1}{3}} \times \sqrt{-1}$

$x = i\sqrt{\frac{1}{3}}$

$x = \frac{\sqrt{3}}{3}\,i$

Therefore, the zeros of $f(x)$ $\frac{\sqrt{3}}{3}\,i, -\frac{\sqrt{3}}{3}\,i,
-\frac{\sqrt{5}}{5}\,i, \frac{\sqrt{5}}{5}\,i$. Since each factor occurs only once, all
the zeros are of multiplicity 1 and the total number of zeros is four.
\end{example}

\begin{example}{Example 3.14}
Factorise the polynomial f$(x) = x^4 - 4x^2 - 21$ completely and find all its zeros.

State the multiplicity of each zero.

\solution
f$(x) = x^4 - 4x^2 - 21$

in this scenario, we can represent $x^2$ by u

$f(x) = (x^2)^2 - 4x^2 - 21$

$f(x) = (u)^2 - 4u - 21$

$u^2 - 4u - 21$

Factorising:

$u^2 - 7u + 3u - 21$

$u(u - 7) + 3(u - 7)$

$(u + 3)(u - 7)$

Substituting $x^2$ back,

$(x^2 + 3)(x^2 - 7)$

To find the zeros, equate the factors to zero:

$(x^2 + 3)(x^2 - 7) = 0$

$x^2 + 3 = 0$ or $x^2 - 7 = 0$

For $x^2 + 3 = 0$

$x^2 = -3$

$x^2 = -3$

$x = \pm\sqrt{-3}$

$x = -\sqrt{-3}$ or $x = \sqrt{-3}$

For $x = -\sqrt{-3}$

$\sqrt{-3} = -\sqrt{3} \times \sqrt{-1}$

$x = -\sqrt{3} \times \sqrt{-1}$

$x = -i\sqrt{3}$

For $x = \sqrt{-3}$

$x = \sqrt{3} \times \sqrt{-1}$

$x = i\sqrt{3}$

For $x^2 - 7$

$x^2 = 7$

$x^2 = 7$

$x = \pm\sqrt{7}$

$x = -\sqrt{7}$ or $x = \sqrt{7}$

For $x = -\sqrt{7}$

$x = -\sqrt{7}$

For $x = \sqrt{7}$

$x = \sqrt{7}$

$x = \sqrt{7}$

Therefore, the zeros of $f(x)$ are $x = i\sqrt{3}$, $x = -i\sqrt{3}$, $x = -\sqrt{7}$,
$x = \sqrt{7}$. Since each factor occurs only once, all the zeros are of multiplicity
1 and the total number of zeros is four.
\end{example}

Note: Multiplicity refers to the number of times a particular zero (or root) appears
in the factorisation of a polynomial. It indicates how many times a specific value
$x = r$ is a solution to the polynomial equation P($x$) $=$ 0.

\noindent\textit{\color{blue}Key Points about Multiplicity:}
\begin{enumerate}
  \item Simple Roots: If a root appears once, it has a multiplicity of 1. For example,
  in the polynomial $(x - 2)$, the root $x = 2$ has multiplicity 1.
  \item Repeated Roots: If a root appears more than once, it has a higher multiplicity.
  For example, in the polynomial $(x - 2)^2$, the root $x = 2$ has a multiplicity of
  2.
  \item Multiplicity and Behaviour: This can be even or odd.
\end{enumerate}
\textit{Odd Multiplicity}: If a root has an odd multiplicity, the graph of the polynomial
crosses the $x$ - axis at that root.

\textit{Even Multiplicity}: If a root has an even multiplicity, the graph of the polynomial
touches the x-axis at that root but does not cross it.

For instance, for the polynomial P(x) $= (x - 1)\,(x - 1)\,(x + 2)$:

The root $x = 1$ has a multiplicity of 2 (it appears twice), so the graph just touches
the x-axis at x $=$ 1 and does not cross it.

The root $x = -2$ has a multiplicity of 1 (it appears once), so the graph crosses the
x-axis at x $=$ -2.
